\section{HNN extensions and indexed parabolics}\label{sec:hnn}


Let $A$ be a group and let $\mathcal P$ be a family of its subgroups,
closed under finite intersections. We say that $\mathcal P$ has the
\emph{coset Helly property} if every finite family of pairwise intersecting
left cosets of members of $\mathcal P$ has nonempty total intersection.
We include $A$ in $\mathcal P$. By Lemma~\ref{F-lem:cosets}, all standard parabolics of a finite-rank flag-transitive coset geometry form such a family.

Fix $C,D\in\mathcal P$ and an isomorphism $\phi:C\longrightarrow D$.
Write
\[
 G=\langle A,t\mid t^{-1}ct=\phi(c)\ (c\in C)\rangle.
\]
A subgroup $P\in\mathcal P$ is called compatible if
$\phi(P\cap C)=P\cap D$. Put $\widehat P=\langle P,t\rangle$.
Britton's lemma embeds the HNN extension of $P$ along these intersections
in $G$ and gives $\widehat P\cap A=P$. These are retained classical
normal-form facts. Unlike an arbitrary family of special subgroups, the
family below keeps its local coset Helly hypothesis.

\begin{theorem}\label{H-thm:helly}
If $\mathcal P$ has the coset Helly property, then the family consisting
of all $P\in\mathcal P$ and all $\widehat P$ for compatible
$P\in\mathcal P$ has the coset Helly property in $G$.
For compatible $P,Q$,
\[
 \widehat P\cap\widehat Q=\widehat{P\cap Q},\qquad
 R\cap\widehat P=R\cap P\quad(R\in\mathcal P).
\]
\end{theorem}
\begin{proof}
Use the Bass--Serre tree $T$ of this HNN extension. Its vertices are
$gA$ and its oriented edges are $gC$, with endpoints $gA$ and $gtA$.
The latter endpoint is independent of the representative because
$ct=t\phi(c)$. At a vertex $xA$, outgoing edge directions are
parametrized by left cosets of $C$ in $A$, and incoming directions by
left cosets of $D$ in $A$.

For compatible $P$, the embedded Bass--Serre tree of $\widehat P$ has
vertices $hA$ and edges $hC$, $h\in\widehat P$. Its injectivity follows
from the same reduced word theorem, since the associated intersections
are exactly $P\cap C$ and $P\cap D$. Denote this connected subtree by
$T_P$. For a base subgroup $R\leq A$ instead use the singleton subtree
$\{A\}$. Attach to a coset $g\widehat P$ the subtree $gT_P$ and to a
coset $gR$ the vertex $gA$. These assignments are independent of the
chosen coset representative.

If a vertex $v=xA$ belongs to such a support subtree, the corresponding
coset meets $xA$. After multiplication by $x^{-1}$ this intersection is
a left coset in $A$ of $P$, or of $R$ in the base-subgroup case.
For example, choose $y\in g\widehat P\cap xA$; the intersection is
$y(\widehat P\cap A)=yP$. An incident edge belongs to $gT_P$ precisely
when this local $P$ coset meets the local $C$ or $D$ coset describing
the direction of that edge. This follows directly from the vertex and
edge descriptions of the embedded tree.

We need one transport observation. Suppose two global cosets meet, and
a vertex $v$ lies in both support subtrees. Choose a common element $z$
of the cosets. Both support subtrees contain $zA$, so they contain the
whole finite path from $zA$ to $v$. Start with the common representative
$z$ at its first vertex. At the next edge, the two local subgroup cosets
intersect each other and each meets the local associated-subgroup coset
of that edge. Local coset Helly supplies a common element of all three.
For an outgoing edge multiply this representative on the right by $t$;
for an incoming edge multiply it by $t^{-1}$. The result lies at the
next vertex and still belongs to both global cosets, because both
subgroups contain $t$. If either subgroup was a base subgroup, its
support was a singleton and the path has length zero, so no multiplication
is needed. Repeating the step gives a common representative in $xA$.

Now take finitely many pairwise intersecting global cosets. Their support
subtrees intersect pairwise. The elementary Helly property for subtrees
gives a common vertex $v=xA$. By the transport observation every pair
of local fibers at $v$ intersects. These fibers are cosets of members
of $\mathcal P$, so local coset Helly supplies an element in all of
them. Multiplying by $x$ proves the first assertion.

For the intersection formula, first observe that $P\cap Q$ is compatible:
an isomorphism commutes with intersections. Thus
$K=\widehat{P\cap Q}$ is contained in both groups. Take
$g\in\widehat P\cap\widehat Q$ and follow the finite path from $A$
to $gA$. Both embedded subtrees contain it. We show inductively that
every vertex on this path is $kA$ for some $k\in K$.
This holds initially with $k=1$. Translate a current vertex $kA$ and
its next edge by $k^{-1}$. If that edge is outgoing, its local coset
$aC$ meets both $P$ and $Q$; if incoming, its local coset $aD$ does.
Since $P\cap Q$ contains $1$, the three local cosets meet pairwise.
The local Helly property supplies $p\in P\cap Q$ in the relevant
coset. The next vertex is therefore $kp tA$ or $kp t^{-1}A$,
respectively, and its representative is again in $K$.
At the end $gA=kA$, and hence
$k^{-1}g\in A\cap\widehat P\cap\widehat Q=P\cap Q\leq K$.
This proves the first equality. The second follows immediately from
$\widehat P\cap A=P$.
\end{proof}

\begin{corollary}\label{H-cor:geometry}
Let $(A,(A_i)_{i\in I})$ be a flag-transitive coset geometry of finite
rank. Let $C,D$ be standard parabolics, and let $\phi:C\to D$ be an
isomorphism. Choose a partition $\mathcal S$ of $I$ such that every
$A_s=\bigcap_{i\in s}A_i$, $s\in\mathcal S$, is compatible, and choose
$F\subseteq I$. In the HNN extension take maximal parabolics
\[
 G_s=\langle A_s,t\rangle\quad(s\in\mathcal S),\qquad
 G_*=A_F.
\]
This is a flag-transitive coset geometry. For a set $J$ of its types,
write $E=\bigcup_{s\in J\cap\mathcal S}s$. Then
\[
 G_J=
 \begin{cases}
   \langle A_E,t\rangle,&*\notin J,\\
   A_{E\cup F},&*\in J.
 \end{cases}
\]
Its chamber stabilizer is $A_I$.
\end{corollary}
\begin{proof}
Take $\mathcal P$ to be all standard parabolics of the input, including
$A$. These have the local coset Helly property explained above. Theorem
\ref{H-thm:helly} applies to all the displayed maximal parabolics. A flag
is a finite family of pairwise intersecting such cosets, so it has a
common representative. It is therefore a translate of a standard flag
and extends to a standard chamber. The intersection formulas follow
by induction from the theorem, with the usual empty-intersection convention.
For the complete type set $E=I$, so the Borel is $A_I$.
\end{proof}

\begin{corollary}\label{H-cor:rc}
Suppose the input in Corollary~\ref{H-cor:geometry} is residually connected.
Let $C=A_J$, $D=A_K$, and suppose a specified bijection
$f:I\setminus J\longrightarrow I\setminus K$ satisfies
\[
 \phi(A_{J\cup X})=A_{K\cup f(X)}\qquad(X\subseteq I\setminus J).
\]
Take $\mathcal S$ to be the connected classes of the arrows $i\to f(i)$
on $I$, and put $F=J\setminus K$. Then the resulting HNN geometry is
residually connected. Its $*$ panels are infinite. Its panel of type
$s\in\mathcal S$ has size
\[
 [A_{I\setminus(s\setminus F)}:A_I].
\]
In particular it is firm when the input is firm.
\end{corollary}
\begin{proof}
Write $H=A_I$ and $R_i=A_{I\setminus\{i\}}$.
The displayed compatibility gives $\phi(H)=H$ and
$\phi(R_i)=R_{f(i)}$ for $i\notin J$.
Each class $s$ satisfies $f(s\setminus J)=s\setminus K$.
Consequently $\phi(A_s\cap C)=A_s\cap D$, so the previous corollary
applies. Its minimal parabolics are
\[
 M_s=A_{I\setminus(s\setminus F)},\qquad M_*=\langle H,t\rangle.
\]

Residual connectivity of the input gives the elementary generation rule
\[
 A_{I\setminus E}=\langle H,R_i\ (i\in E)\rangle
 \quad(E\subseteq I).
\]
This is Lemma~\ref{O-lem:parabolic-join}, including the empty support; the Borel is not dropped.

Let $W$ be a subset of the new type set $\mathcal S\cup\{*\}$, and
put $E=\bigcup_{s\in W\cap\mathcal S}s$. If $*\notin W$, then
\[
 \langle H,M_s\ (s\in W)\rangle
 =A_{I\setminus(E\setminus F)}=G_{(\mathcal S\cup\{*\})\setminus W}.
\]
If $*\in W$, this generated group also contains $t$.
Every $i\in E\cap F$ is in the range but not the domain of $f$.
It therefore has a predecessor $j\in E\setminus F$ with $f(j)=i$.
Conjugation by $t$ brings $R_i$ into the group from $R_j$.
Thus in this case the generated group is
$\langle H,t,R_i\ (i\in E)\rangle=\langle A_{I\setminus E},t\rangle$,
again exactly $G_{(\mathcal S\cup\{*\})\setminus W}$.

We have proved the same generation rule for every new standard parabolic.
If at least two types remain, each of its minimal generators belongs to
one of the parabolics obtained by deleting a different remaining type.
Those parabolics therefore generate it. The residual generation criterion
proves residual connectivity.

Since $\phi(H)=H$, the subgroup $\langle H,t\rangle$ is an extension
of $H$ by the infinite cyclic group generated by the stable letter.
Its index over $H$ is infinite. The other panel indices are the stated
indices of $M_s$ over $H$. Every class has a point outside $F$: a point
in $F$ has a predecessor outside $F$. If the input is firm, the group
$M_s$ therefore contains some $R_i$ strictly larger than $H$.
All these panels have at least two elements, as required.
\end{proof}


\subsection{Firmness recovers the missing indices}


Let $\alpha=(A,(A_i)_{i\in I})$ be a coset geometry of finite rank.
Write $A_J=\bigcap_{j\in J}A_j$, $A_\varnothing=A$,
$H=A_I$, and $R_i=A_{I\setminus\{i\}}$.
The following criterion turns an isomorphism of unlabelled parabolic lattices into the indexed compatibility needed above. Firmness is an actual hypothesis.

\begin{lemma}\label{I-lem:boolean}
Suppose $R_i>H$ for every $i\in I$. Then
\[
 A_J\leq A_K\quad\Longleftrightarrow\quad K\subseteq J.
\]
Consequently the lattice of standard parabolics, ordered by inclusion,
is the reverse Boolean lattice on $I$.
\end{lemma}
\begin{proof}
If $K\subseteq J$, the inclusion follows from the definition of an
intersection. Conversely, suppose $k\in K\setminus J$. Then
$R_k\leq A_J$, because every member of $J$ differs from $k$.
If $A_J\leq A_K$, it follows that $R_k\leq A_k$, whence
$R_k=R_k\cap A_k=H$, a contradiction. This proves the equivalence,
including injectivity of $J\mapsto A_J$.
\end{proof}

\begin{corollary}\label{I-cor:firm}
Suppose $\alpha$ is flag-transitive and firm. Let
$\phi:A_J\longrightarrow A_K$ be an isomorphism inducing an isomorphism
of the unlabelled lattices of standard parabolics contained in its domain
and codomain. Then there is a unique bijection
$f:I\setminus J\longrightarrow I\setminus K$ such that
\[
 \phi(A_{J\cup X})=A_{K\cup f(X)}
 \qquad(X\subseteq I\setminus J).
\]
Thus the indexed hypothesis of Corollary~\ref{H-cor:rc} follows from the source's
admissibility condition for firm flag-transitive inputs.
\end{corollary}
\begin{proof}
Flag transitivity identifies the chambers on a standard panel of missing
type $i$ with $R_i/H$. Firmness therefore implies $[R_i:H]\geq2$.
Lemma~\ref{I-lem:boolean} applies. The standard parabolics contained in
$A_J$ are precisely the $A_{J\cup X}$, and their coatoms are
$A_{J\cup\{i\}}$ for $i\notin J$. The isomorphism permutes coatoms
bijectively, giving $f$. Since it preserves intersections, intersecting
these coatom equalities proves the formula for every $X$, including
$X=\varnothing$.

For completeness, form the equivalence classes $s$ of the arrows
$i\mapsto f(i)$ on $I$. Bijection of the domain and range gives
$f(s\setminus J)=s\setminus K$. Hence
\[
 \phi(A_s\cap A_J)
 =\bigcap_{i\in s\setminus J}\phi(A_{J\cup\{i\}})
 =\bigcap_{j\in s\setminus K}A_{K\cup\{j\}}
 =A_s\cap A_K.
\]
Both displayed intersections of images are taken inside $A_K$; if
their index sets are empty they mean $A_K$, in agreement with
$\phi(A_J)=A_K$. Every $A_s$ is therefore compatible. These are precisely the
subgroups required by the local coset Helly theorem. Its flag-transitivity
conclusion follows; its residual-connectivity conclusion also requires
residual connectivity of the input, as stated there.
\end{proof}


\subsection{Residually connected inputs with universal types}
Let $\alpha=(A,(A_i)_{i\in I})$ be a finite-rank flag-transitive,
residually connected coset geometry. Write $A_J=\bigcap_{i\in J}A_i$,
$H=A_I$, and $R_i=A_{I\setminus\{i\}}$.
The Borel-preserving residual generation formula is
\begin{equation}\label{R-eq:gen}
 A_J=\langle H,R_i:i\notin J\rangle.
\end{equation}
This is the same Borel-preserving formula proved in Lemma~\ref{O-lem:parabolic-join}.

\begin{lemma}\label{R-lem:active}
Put $I_0=\{i:A_i=A\}$ and $I_+=I\setminus I_0$.
Then $A_i=A$ if and only if $R_i=H$. Deleting the universal types
$I_0$ leaves a firm flag-transitive, residually connected geometry
on $I_+$ with the same unlabelled family of standard parabolics and
the same Borel $H$. The original incidence geometry is its join
with one singleton for each type in $I_0$.
\end{lemma}
\begin{proof}
If $A_i=A$, then $H=A_i\cap R_i=R_i$. Conversely, if $R_i=H$,
formula (\ref{R-eq:gen}) for $J=\varnothing$ and $J=\{i\}$ gives
$A=A_i$. A type with maximal parabolic $A$ has exactly one element,
incident with every other element, proving the join description.
Dropping these factors leaves every standard intersection unchanged.
Flags still have common representatives. Every residue on at least
two active types is a residue of the original geometry after fixing
the deleted universal elements, so it remains connected. Each active
panel has index $[R_i:H]>1$, proving firmness. The empty active
geometry is interpreted as a single empty chamber.
\end{proof}

For later comparison we also make the firm branch explicit without
assuming residual connectedness. For a firm flag-transitive input
and a source-admissible $\phi:A_J\to A_K$,
Lemma~\ref{I-lem:boolean} gives the indexed partial
bijection required by Theorem~\ref{H-thm:helly}. For its classes
$s$ and $F=J\setminus K$, that theorem's intersection formula gives
the Borel $H$, the stable minimal parabolic $\langle H,t\rangle$,
and the other minimal parabolics
\[
 M_s=A_{I\setminus(s\setminus F)}.
\]
The first has infinite index over $H$, since the stable letter has
infinite image in $\ZZ$ and normalizes $H$. Every class has an
element outside $F$: a point in $F$ is in the range but not the
domain of the partial bijection, and its predecessor is outside
$J$, hence outside $F$. Consequently $M_s$ contains some
$R_i>H$. All panels are firm. This proves the full firm-input branch
of source Theorem 3.35, including part (c), independently of the
residual-connectivity assertion below.

Suppose $C=A_J,D=A_K$, and $\phi:C\to D$ satisfies the source's
Definition 3.30: it induces an isomorphism of the unlabelled families
of standard parabolics lying in these subgroups. The source forms
classes of the arrows $i\to j$, with $i\notin J,j\notin K$, for
which $\phi(A_{J\cup\{i\}})=A_{K\cup\{j\}}$. It uses
$F=J\setminus K$, the maximal subgroups $\langle A_s,t\rangle$
for these connected classes $s$, and the stable-type subgroup $A_F$
in the HNN extension $\langle A,t\mid t^{-1}ct=\phi(c)\rangle$.

\begin{theorem}\label{R-thm:rc}
With these literal choices, every flag-transitive, residually connected
input gives a flag-transitive, residually connected output. Thus the
residually connected branch of Theorem 3.35 remains valid without
firmness. Active output panels have the indices in the indexed HNN theorem; each additional universal class has one element
in its panel.
\end{theorem}
\begin{proof}
Apply Lemma~\ref{R-lem:active}. On the active types the firm geometry
has the reverse Boolean parabolic lattice proved in Lemma~\ref{I-lem:boolean}. Thus the isomorphism $\phi$ gives a
bijection
\[
 f:I_+\setminus J\longrightarrow I_+\setminus K
\]
with
$\phi(A_{J\cup X})=A_{K\cup f(X)}$ for every
$X\subseteq I_+\setminus J$. Here intersections with $J,K$
can ignore their inactive entries. This follows because the coatoms
of the Boolean interval inside $A_J$ are exactly the subgroups
$A_{J\cup\{i\}}$ for its missing active types; an isomorphism
permutes those coatoms and commutes with all their intersections.
The minimal subgroup $H$ is also preserved.

For an active $i\notin J$, the subgroup $A_{J\cup\{i\}}$ is
proper in $A_J$, whereas for an inactive $i$ it equals $A_J$.
The same distinction holds on the $K$ side. Hence a source arrow
never connects an active type to an inactive one. On active types
the connected classes are exactly those of the partial bijection
$f$, with isolated active types included. The stable parabolic is
$A_F=A_{F\cap I_+}$. Thus all active-class parabolics and the
stable parabolic coincide exactly with the indexed HNN
construction of Corollary~\ref{H-cor:rc} on the firm active geometry. Its flag transitivity and
residual connectedness follow from that theorem.

An inactive class $s$ consists only of types with $A_i=A$, so
$A_s=A$ and $\langle A_s,t\rangle$ is the entire HNN group.
Its coset type is a singleton universal element. Adding any number
of these types preserves flag transitivity and residual connectedness:
a residue with at least two active types uses the active conclusion,
a residue with one active type and at least one universal type has
a connected star, and a residue consisting of at least two universal
types is a complete graph. The same observation covers the case
$I_+=\varnothing$. In that case $C=D=A$ and the stable-type cosets
of $A$ in $A\rtimes\ZZ$ are the infinite rank-one part; all other
types are universal. This proves the assertion and panel description.
\end{proof}
